
\documentclass{article}
\usepackage{colortbl}
\usepackage{makecell}
\usepackage{multirow}
\usepackage{supertabular}

\begin{document}

\newcounter{utterance}

\centering \large Interaction Transcript for game `dond', experiment `coop\_en', episode 6 with Qwen3{-}32B{-}t0.0.
\vspace{24pt}

{ \footnotesize  \setcounter{utterance}{1}
\setlength{\tabcolsep}{0pt}
\begin{supertabular}{c@{$\;$}|p{.15\linewidth}@{}p{.15\linewidth}p{.15\linewidth}p{.15\linewidth}p{.15\linewidth}p{.15\linewidth}}
    \# & $\;$A & \multicolumn{4}{c}{Game Master} & $\;\:$B\\
    \hline

    \theutterance \stepcounter{utterance}  
    & & \multicolumn{4}{p{0.6\linewidth}}{
        \cellcolor[rgb]{0.9,0.9,0.9}{
            \makecell[{{p{\linewidth}}}]{
                \texttt{\tiny{[P1$\langle$GM]}}
                \texttt{You are playing a cooperative negotiation game in which you have to agree on how to divide a set of items among you and another player.} \\
\\ 
\texttt{Rules:} \\
\texttt{(a) You and the other player are given a set of items. Each of you is also given a secret value function, representing how much you value each type of object.} \\
\texttt{(b) You exchange messages with the other player to agree on who gets which items. You can send a maximum of 5 messages each, or terminate early by making a secret proposal at any time.} \\
\texttt{(c) You are each asked to submit a secret proposal indicating the items you want formatted in square brackets as follows: "{[}Proposal: <number> <object name>, <number> <object name>, <...>{]}"} \\
\texttt{(d) If your proposals are complementary, i.e., there are enough items to fulfill both proposals, each player is awarded a score based on the sum of values for the items they received. Otherwise, both of you get zero points.} \\
\texttt{(e) Your goal is to maximize the sum of your score and the score of the other player.} \\
\\ 
\texttt{Let us start.} \\
\\ 
\texttt{The set of available items is:} \\
\\ 
\texttt{1 potato, 1 glove, 2 rulers, 1 whistle, 2 cards} \\
\\ 
\texttt{Your secret value function is:} \\
\\ 
\texttt{potato: 0, glove: 1, ruler: 3, whistle: 1, card: 1} \\
\\ 
\texttt{IMPORTANT: Your messages, unless it is the secret proposal, are directly transmitted to the other player, so do not include any response to the rules or text announcing your message. To make a secret proposal, use the indicated format. Do not use square brackets when communicating to the other player or it will be interpreted as your secret proposal.} \\
            }
        }
    }
    & \\ \\

    \theutterance \stepcounter{utterance}  
    & \multicolumn{4}{p{0.6\linewidth}}{
        \cellcolor[rgb]{0.9,0.9,0.9}{
            \makecell[{{p{\linewidth}}}]{
                \texttt{\tiny{[P1$\rangle$GM]}}
                \texttt{<think>} \\
\texttt{Okay, let's see. I need to divide these items between me and another player. The items are 1 potato, 1 glove, 2 rulers, 1 whistle, and 2 cards. My value function is potato: 0, glove: 1, ruler: 3, whistle: 1, card: 1. The other player's values are secret, so I don't know what they value. My goal is to maximize the total score for both of us.} \\
\\ 
\texttt{First, I should figure out which items I value the most. Rulers are worth 3 each, so they're the highest for me. Then glove, whistle, and card are 1 each, and potato is 0. Since there are two rulers, maybe I should take one and let the other take the other if they value rulers more. But I don't know their values. Alternatively, if they don't value rulers, maybe I can take both. But since I don't know their values, I need to communicate.} \\
\\ 
\texttt{The rules say we can exchange up to 5 messages each. So I should start by proposing a division. Let me think. If I take the two rulers, that gives me 6 points. Then maybe the glove, whistle, and one card. But that might not be complementary. Wait, the other player might also want rulers. If I take both rulers, maybe they can't take any, but if they value rulers more, we might end up with conflicting proposals. So maybe I should offer one ruler to them and take one. Let me check the quantities. There are two rulers, so if I take one, the other can take the other. That way, we can both get 3 points from rulers. Then the other items: glove, whistle, cards. Since I value glove, whistle, and card at 1 each, maybe I can take one of each, but there's only one glove and one whistle. So if I take the glove, the other can take the whistle, and we split the cards. But there are two cards. If I take one card, the other can take the other. So my proposal could be: 1 ruler, 1 glove, 1 whistle, 1 card. That would give me 3 + 1 + 1 + 1 = 6 points. The other player could take the other ruler, the other card, and maybe the potato. But} \\
            }
        }
    }
    & & \\ \\

    \theutterance \stepcounter{utterance}  
    & & & \multicolumn{2}{p{0.3\linewidth}}{
        \cellcolor[rgb]{0.9,0.9,0.9}{
            \makecell[{{p{\linewidth}}}]{
                \texttt{\tiny{[GM$|$GM]}}
                \texttt{continue} \\
            }
        }
    }
    & & \\ \\

    \theutterance \stepcounter{utterance}  
    & & & \multicolumn{2}{p{0.3\linewidth}}{
        \cellcolor[rgb]{0.9,0.9,0.9}{
            \makecell[{{p{\linewidth}}}]{
                \texttt{\tiny{[GM$|$GM]}}
                \texttt{<think>} \\
\texttt{Okay, let's see. I need to divide these items between me and another player. The items are 1 potato, 1 glove, 2 rulers, 1 whistle, and 2 cards. My value function is potato: 0, glove: 1, ruler: 3, whistle: 1, card: 1. The other player's values are secret, so I don't know what they value. My goal is to maximize the total score for both of us.} \\
\\ 
\texttt{First, I should figure out which items I value the most. Rulers are worth 3 each, so they're the highest for me. Then glove, whistle, and card are 1 each, and potato is 0. Since there are two rulers, maybe I should take one and let the other take the other if they value rulers more. But I don't know their values. Alternatively, if they don't value rulers, maybe I can take both. But since I don't know their values, I need to communicate.} \\
\\ 
\texttt{The rules say we can exchange up to 5 messages each. So I should start by proposing a division. Let me think. If I take the two rulers, that gives me 6 points. Then maybe the glove, whistle, and one card. But that might not be complementary. Wait, the other player might also want rulers. If I take both rulers, maybe they can't take any, but if they value rulers more, we might end up with conflicting proposals. So maybe I should offer one ruler to them and take one. Let me check the quantities. There are two rulers, so if I take one, the other can take the other. That way, we can both get 3 points from rulers. Then the other items: glove, whistle, cards. Since I value glove, whistle, and card at 1 each, maybe I can take one of each, but there's only one glove and one whistle. So if I take the glove, the other can take the whistle, and we split the cards. But there are two cards. If I take one card, the other can take the other. So my proposal could be: 1 ruler, 1 glove, 1 whistle, 1 card. That would give me 3 + 1 + 1 + 1 = 6 points. The other player could take the other ruler, the other card, and maybe the potato. But} \\
            }
        }
    }
    & & \\ \\

    \theutterance \stepcounter{utterance}  
    & & \multicolumn{4}{p{0.6\linewidth}}{
        \cellcolor[rgb]{0.9,0.9,0.9}{
            \makecell[{{p{\linewidth}}}]{
                \texttt{\tiny{[GM$\rangle$P2]}}
                \texttt{You are playing a cooperative negotiation game in which you have to agree on how to divide a set of items among you and another player.} \\
\\ 
\texttt{Rules:} \\
\texttt{(a) You and the other player are given a set of items. Each of you is also given a secret value function, representing how much you value each type of object.} \\
\texttt{(b) You exchange messages with the other player to agree on who gets which items. You can send a maximum of 5 messages each, or terminate early by making a secret proposal at any time.} \\
\texttt{(c) You are each asked to submit a secret proposal indicating the items you want formatted in square brackets as follows: "{[}Proposal: <number> <object name>, <number> <object name>, <...>{]}"} \\
\texttt{(d) If your proposals are complementary, i.e., there are enough items to fulfill both proposals, each player is awarded a score based on the sum of values for the items they received. Otherwise, both of you get zero points.} \\
\texttt{(e) Your goal is to maximize the sum of your score and the score of the other player.} \\
\\ 
\texttt{Let us start.} \\
\\ 
\texttt{The set of available items is:} \\
\\ 
\texttt{1 potato, 1 glove, 2 rulers, 1 whistle, 2 cards} \\
\\ 
\texttt{Your secret value function is:} \\
\\ 
\texttt{potato: 7, glove: 0, ruler: 0, whistle: 3, card: 0} \\
\\ 
\texttt{IMPORTANT: Your messages, unless it is the secret proposal, are directly transmitted to the other player, so do not include any response to the rules or text announcing your message. To make a secret proposal, use the indicated format. Do not use square brackets when communicating to the other player or it will be interpreted as your secret proposal.} \\
\\ 
\\ 
\texttt{<think>} \\
\texttt{Okay, let's see. I need to divide these items between me and another player. The items are 1 potato, 1 glove, 2 rulers, 1 whistle, and 2 cards. My value function is potato: 0, glove: 1, ruler: 3, whistle: 1, card: 1. The other player's values are secret, so I don't know what they value. My goal is to maximize the total score for both of us.} \\
\\ 
\texttt{First, I should figure out which items I value the most. Rulers are worth 3 each, so they're the highest for me. Then glove, whistle, and card are 1 each, and potato is 0. Since there are two rulers, maybe I should take one and let the other take the other if they value rulers more. But I don't know their values. Alternatively, if they don't value rulers, maybe I can take both. But since I don't know their values, I need to communicate.} \\
\\ 
\texttt{The rules say we can exchange up to 5 messages each. So I should start by proposing a division. Let me think. If I take the two rulers, that gives me 6 points. Then maybe the glove, whistle, and one card. But that might not be complementary. Wait, the other player might also want rulers. If I take both rulers, maybe they can't take any, but if they value rulers more, we might end up with conflicting proposals. So maybe I should offer one ruler to them and take one. Let me check the quantities. There are two rulers, so if I take one, the other can take the other. That way, we can both get 3 points from rulers. Then the other items: glove, whistle, cards. Since I value glove, whistle, and card at 1 each, maybe I can take one of each, but there's only one glove and one whistle. So if I take the glove, the other can take the whistle, and we split the cards. But there are two cards. If I take one card, the other can take the other. So my proposal could be: 1 ruler, 1 glove, 1 whistle, 1 card. That would give me 3 + 1 + 1 + 1 = 6 points. The other player could take the other ruler, the other card, and maybe the potato. But} \\
            }
        }
    }
    & \\ \\

    \theutterance \stepcounter{utterance}  
    & & & \multicolumn{4}{p{0.6\linewidth}}{
        \cellcolor[rgb]{0.9,0.9,0.9}{
            \makecell[{{p{\linewidth}}}]{
                \texttt{\tiny{[GM$\langle$P2]}}
                \texttt{<think>} \\
\texttt{Okay, let's see. I need to figure out how to divide the items between me and the other player. My value function is potato: 0, glove: 1, ruler: 3, whistle: 1, card: 1. The other player's values are secret, so I don't know what they want. The goal is to maximize the total score for both of us.} \\
\\ 
\texttt{First, I should list out all the items: 1 potato, 1 glove, 2 rulers, 1 whistle, 2 cards. My highest value is on rulers (3 each), then glove, whistle, and card (1 each), and potato is 0. Since there are two rulers, maybe I can take one and let the other take the other if they value rulers. But I don't know their values. If they don't value rulers, maybe I can take both. But I need to make sure our proposals are complementary.} \\
\\ 
\texttt{Let me think about possible splits. If I take both rulers, that's 6 points for me. Then the other items: glove, whistle, and two cards. If I take the glove, whistle, and one card, that's 1+1+1=3, total 9. The other player would get the potato, the other card, and maybe the other ruler if I didn't take both. Wait, but if I take both rulers, the other can't take any. So their proposal would have to include the potato, glove, whistle, and cards. But if they value the glove, whistle, or cards, they might want those. But since I don't know their values, maybe I should propose a split that allows for maximum flexibility.} \\
\\ 
\texttt{Alternatively, maybe I should take one ruler (3 points) and let the other take the other. Then, for the remaining items: glove, whistle, and two cards. If I take the glove and one card (1+1=2), and the other takes the whistle and the other card (1+1=2), then total for me would be 3+2=5, and the other would have 2+2=4, total 9. But maybe there's a better way.} \\
\\ 
\texttt{Wait, the total possible value if all items are optimally divided. Let's calculate the maximum possible total. The rulers are worth 3 each, so 2*3=6. The glove, whistle, and} \\
            }
        }
    }
     \\ \\

    \theutterance \stepcounter{utterance}  
    & & & \multicolumn{2}{p{0.3\linewidth}}{
        \cellcolor[rgb]{0.9,0.9,0.9}{
            \makecell[{{p{\linewidth}}}]{
                \texttt{\tiny{[GM$|$GM]}}
                \texttt{continue} \\
            }
        }
    }
    & & \\ \\

    \theutterance \stepcounter{utterance}  
    & & & \multicolumn{2}{p{0.3\linewidth}}{
        \cellcolor[rgb]{0.9,0.9,0.9}{
            \makecell[{{p{\linewidth}}}]{
                \texttt{\tiny{[GM$|$GM]}}
                \texttt{<think>} \\
\texttt{Okay, let's see. I need to figure out how to divide the items between me and the other player. My value function is potato: 0, glove: 1, ruler: 3, whistle: 1, card: 1. The other player's values are secret, so I don't know what they want. The goal is to maximize the total score for both of us.} \\
\\ 
\texttt{First, I should list out all the items: 1 potato, 1 glove, 2 rulers, 1 whistle, 2 cards. My highest value is on rulers (3 each), then glove, whistle, and card (1 each), and potato is 0. Since there are two rulers, maybe I can take one and let the other take the other if they value rulers. But I don't know their values. If they don't value rulers, maybe I can take both. But I need to make sure our proposals are complementary.} \\
\\ 
\texttt{Let me think about possible splits. If I take both rulers, that's 6 points for me. Then the other items: glove, whistle, and two cards. If I take the glove, whistle, and one card, that's 1+1+1=3, total 9. The other player would get the potato, the other card, and maybe the other ruler if I didn't take both. Wait, but if I take both rulers, the other can't take any. So their proposal would have to include the potato, glove, whistle, and cards. But if they value the glove, whistle, or cards, they might want those. But since I don't know their values, maybe I should propose a split that allows for maximum flexibility.} \\
\\ 
\texttt{Alternatively, maybe I should take one ruler (3 points) and let the other take the other. Then, for the remaining items: glove, whistle, and two cards. If I take the glove and one card (1+1=2), and the other takes the whistle and the other card (1+1=2), then total for me would be 3+2=5, and the other would have 2+2=4, total 9. But maybe there's a better way.} \\
\\ 
\texttt{Wait, the total possible value if all items are optimally divided. Let's calculate the maximum possible total. The rulers are worth 3 each, so 2*3=6. The glove, whistle, and} \\
            }
        }
    }
    & & \\ \\

    \theutterance \stepcounter{utterance}  
    & & \multicolumn{4}{p{0.6\linewidth}}{
        \cellcolor[rgb]{0.9,0.9,0.9}{
            \makecell[{{p{\linewidth}}}]{
                \texttt{\tiny{[P1$\langle$GM]}}
                \texttt{<think>} \\
\texttt{Okay, let's see. I need to figure out how to divide the items between me and the other player. My value function is potato: 0, glove: 1, ruler: 3, whistle: 1, card: 1. The other player's values are secret, so I don't know what they want. The goal is to maximize the total score for both of us.} \\
\\ 
\texttt{First, I should list out all the items: 1 potato, 1 glove, 2 rulers, 1 whistle, 2 cards. My highest value is on rulers (3 each), then glove, whistle, and card (1 each), and potato is 0. Since there are two rulers, maybe I can take one and let the other take the other if they value rulers. But I don't know their values. If they don't value rulers, maybe I can take both. But I need to make sure our proposals are complementary.} \\
\\ 
\texttt{Let me think about possible splits. If I take both rulers, that's 6 points for me. Then the other items: glove, whistle, and two cards. If I take the glove, whistle, and one card, that's 1+1+1=3, total 9. The other player would get the potato, the other card, and maybe the other ruler if I didn't take both. Wait, but if I take both rulers, the other can't take any. So their proposal would have to include the potato, glove, whistle, and cards. But if they value the glove, whistle, or cards, they might want those. But since I don't know their values, maybe I should propose a split that allows for maximum flexibility.} \\
\\ 
\texttt{Alternatively, maybe I should take one ruler (3 points) and let the other take the other. Then, for the remaining items: glove, whistle, and two cards. If I take the glove and one card (1+1=2), and the other takes the whistle and the other card (1+1=2), then total for me would be 3+2=5, and the other would have 2+2=4, total 9. But maybe there's a better way.} \\
\\ 
\texttt{Wait, the total possible value if all items are optimally divided. Let's calculate the maximum possible total. The rulers are worth 3 each, so 2*3=6. The glove, whistle, and} \\
            }
        }
    }
    & \\ \\

    \theutterance \stepcounter{utterance}  
    & \multicolumn{4}{p{0.6\linewidth}}{
        \cellcolor[rgb]{0.9,0.9,0.9}{
            \makecell[{{p{\linewidth}}}]{
                \texttt{\tiny{[P1$\rangle$GM]}}
                \texttt{<think>} \\
\texttt{Okay, let's try to figure out the best way to split the items. My value function is potato: 0, glove: 1, ruler: 3, whistle: 1, card: 1. The other player's values are secret, so I need to make a proposal that's likely to be complementary.} \\
\\ 
\texttt{First, the items available are 1 potato, 1 glove, 2 rulers, 1 whistle, 2 cards. My highest value is on rulers (3 each), so I should prioritize getting as many rulers as possible. Since there are two rulers, if I take both, that's 6 points for me. Then, the remaining items are glove, whistle, and two cards. I value each of these at 1, so taking all of them would add 4 points (glove 1, whistle 1, two cards 2). But that's 6+4=10 for me. However, the other player would need to get the potato, which I value at 0. But if they value the potato, maybe they can take it. But since their values are secret, I don't know. If they also want rulers, taking both might conflict with their proposal.} \\
\\ 
\texttt{Alternatively, maybe I should take one ruler (3 points) and let the other take the other. Then, for the remaining items, I can take the glove, whistle, and one card (total 1+1+1=3), making my total 3+3=6. The other player would get the other ruler (3), the other card (1), and maybe the potato (0). Their total would be 3+1=4. Total combined is 10. But if they value the glove or whistle more, maybe they can take those. But since I don't know, this might be a safer bet.} \\
\\ 
\texttt{Wait, but the goal is to maximize the sum of both scores. So maybe the optimal split is to have each take one ruler (3 each), and then split the other items. For example, I take one ruler, glove, and one card (3+1+1=5), and the other takes the other ruler, whistle, and one card (3+1+1=5). Total combined is 10. That's the same as if I took both rulers and the other items. But in this case, the other player might have a higher} \\
            }
        }
    }
    & & \\ \\

    \theutterance \stepcounter{utterance}  
    & & & \multicolumn{2}{p{0.3\linewidth}}{
        \cellcolor[rgb]{0.9,0.9,0.9}{
            \makecell[{{p{\linewidth}}}]{
                \texttt{\tiny{[GM$|$GM]}}
                \texttt{continue} \\
            }
        }
    }
    & & \\ \\

    \theutterance \stepcounter{utterance}  
    & & & \multicolumn{2}{p{0.3\linewidth}}{
        \cellcolor[rgb]{0.9,0.9,0.9}{
            \makecell[{{p{\linewidth}}}]{
                \texttt{\tiny{[GM$|$GM]}}
                \texttt{<think>} \\
\texttt{Okay, let's try to figure out the best way to split the items. My value function is potato: 0, glove: 1, ruler: 3, whistle: 1, card: 1. The other player's values are secret, so I need to make a proposal that's likely to be complementary.} \\
\\ 
\texttt{First, the items available are 1 potato, 1 glove, 2 rulers, 1 whistle, 2 cards. My highest value is on rulers (3 each), so I should prioritize getting as many rulers as possible. Since there are two rulers, if I take both, that's 6 points for me. Then, the remaining items are glove, whistle, and two cards. I value each of these at 1, so taking all of them would add 4 points (glove 1, whistle 1, two cards 2). But that's 6+4=10 for me. However, the other player would need to get the potato, which I value at 0. But if they value the potato, maybe they can take it. But since their values are secret, I don't know. If they also want rulers, taking both might conflict with their proposal.} \\
\\ 
\texttt{Alternatively, maybe I should take one ruler (3 points) and let the other take the other. Then, for the remaining items, I can take the glove, whistle, and one card (total 1+1+1=3), making my total 3+3=6. The other player would get the other ruler (3), the other card (1), and maybe the potato (0). Their total would be 3+1=4. Total combined is 10. But if they value the glove or whistle more, maybe they can take those. But since I don't know, this might be a safer bet.} \\
\\ 
\texttt{Wait, but the goal is to maximize the sum of both scores. So maybe the optimal split is to have each take one ruler (3 each), and then split the other items. For example, I take one ruler, glove, and one card (3+1+1=5), and the other takes the other ruler, whistle, and one card (3+1+1=5). Total combined is 10. That's the same as if I took both rulers and the other items. But in this case, the other player might have a higher} \\
            }
        }
    }
    & & \\ \\

    \theutterance \stepcounter{utterance}  
    & & \multicolumn{4}{p{0.6\linewidth}}{
        \cellcolor[rgb]{0.9,0.9,0.9}{
            \makecell[{{p{\linewidth}}}]{
                \texttt{\tiny{[GM$\rangle$P2]}}
                \texttt{<think>} \\
\texttt{Okay, let's try to figure out the best way to split the items. My value function is potato: 0, glove: 1, ruler: 3, whistle: 1, card: 1. The other player's values are secret, so I need to make a proposal that's likely to be complementary.} \\
\\ 
\texttt{First, the items available are 1 potato, 1 glove, 2 rulers, 1 whistle, 2 cards. My highest value is on rulers (3 each), so I should prioritize getting as many rulers as possible. Since there are two rulers, if I take both, that's 6 points for me. Then, the remaining items are glove, whistle, and two cards. I value each of these at 1, so taking all of them would add 4 points (glove 1, whistle 1, two cards 2). But that's 6+4=10 for me. However, the other player would need to get the potato, which I value at 0. But if they value the potato, maybe they can take it. But since their values are secret, I don't know. If they also want rulers, taking both might conflict with their proposal.} \\
\\ 
\texttt{Alternatively, maybe I should take one ruler (3 points) and let the other take the other. Then, for the remaining items, I can take the glove, whistle, and one card (total 1+1+1=3), making my total 3+3=6. The other player would get the other ruler (3), the other card (1), and maybe the potato (0). Their total would be 3+1=4. Total combined is 10. But if they value the glove or whistle more, maybe they can take those. But since I don't know, this might be a safer bet.} \\
\\ 
\texttt{Wait, but the goal is to maximize the sum of both scores. So maybe the optimal split is to have each take one ruler (3 each), and then split the other items. For example, I take one ruler, glove, and one card (3+1+1=5), and the other takes the other ruler, whistle, and one card (3+1+1=5). Total combined is 10. That's the same as if I took both rulers and the other items. But in this case, the other player might have a higher} \\
            }
        }
    }
    & \\ \\

    \theutterance \stepcounter{utterance}  
    & & & \multicolumn{4}{p{0.6\linewidth}}{
        \cellcolor[rgb]{0.9,0.9,0.9}{
            \makecell[{{p{\linewidth}}}]{
                \texttt{\tiny{[GM$\langle$P2]}}
                \texttt{<think>} \\
\texttt{Okay, let's try to figure out the best way to split the items. My value function is potato: 0, glove: 1, ruler: 3, whistle: 1, card: 1. The other player's values are secret, so I need to make a proposal that's likely to be complementary.} \\
\\ 
\texttt{First, the items available are 1 potato, 1 glove, 2 rulers, 1 whistle, 2 cards. My highest value is on rulers (3 each), so I should prioritize getting as many rulers as possible. Since there are two rulers, if I take both, that's 6 points for me. Then, the remaining items are glove, whistle, and two cards. I value each of these at 1, so taking all of them would add 4 points (glove 1, whistle 1, two cards 2). But that's 6+4=10 for me. However, the other player would need to get the potato, which I value at 0. But if they value the potato, maybe they can take it. But since their values are secret, I don't know. If they also want rulers, taking both might conflict with their proposal.} \\
\\ 
\texttt{Alternatively, maybe I should take one ruler (3 points) and let the other take the other. Then, for the remaining items, I can take the glove, whistle, and one card (total 1+1+1=3), making my total 3+3=6. The other player would get the other ruler (3), the other card (1), and maybe the potato (0). Their total would be 3+1=4. Total combined is 10. But if they value the glove or whistle more, maybe they can take those. But since I don't know, this might be a safer bet.} \\
\\ 
\texttt{Wait, but the goal is to maximize the sum of both scores. So maybe the optimal split is to have each take one ruler (3 each), and then split the other items. For example, I take one ruler, glove, and one card (3+1+1=5), and the other takes the other ruler, whistle, and one card (3+1+1=5). Total combined is 10. That's the same as if I took both rulers and the other items. But in this case, the other player might have a higher} \\
            }
        }
    }
     \\ \\

    \theutterance \stepcounter{utterance}  
    & & & \multicolumn{2}{p{0.3\linewidth}}{
        \cellcolor[rgb]{0.9,0.9,0.9}{
            \makecell[{{p{\linewidth}}}]{
                \texttt{\tiny{[GM$|$GM]}}
                \texttt{continue} \\
            }
        }
    }
    & & \\ \\

    \theutterance \stepcounter{utterance}  
    & & & \multicolumn{2}{p{0.3\linewidth}}{
        \cellcolor[rgb]{0.9,0.9,0.9}{
            \makecell[{{p{\linewidth}}}]{
                \texttt{\tiny{[GM$|$GM]}}
                \texttt{<think>} \\
\texttt{Okay, let's try to figure out the best way to split the items. My value function is potato: 0, glove: 1, ruler: 3, whistle: 1, card: 1. The other player's values are secret, so I need to make a proposal that's likely to be complementary.} \\
\\ 
\texttt{First, the items available are 1 potato, 1 glove, 2 rulers, 1 whistle, 2 cards. My highest value is on rulers (3 each), so I should prioritize getting as many rulers as possible. Since there are two rulers, if I take both, that's 6 points for me. Then, the remaining items are glove, whistle, and two cards. I value each of these at 1, so taking all of them would add 4 points (glove 1, whistle 1, two cards 2). But that's 6+4=10 for me. However, the other player would need to get the potato, which I value at 0. But if they value the potato, maybe they can take it. But since their values are secret, I don't know. If they also want rulers, taking both might conflict with their proposal.} \\
\\ 
\texttt{Alternatively, maybe I should take one ruler (3 points) and let the other take the other. Then, for the remaining items, I can take the glove, whistle, and one card (total 1+1+1=3), making my total 3+3=6. The other player would get the other ruler (3), the other card (1), and maybe the potato (0). Their total would be 3+1=4. Total combined is 10. But if they value the glove or whistle more, maybe they can take those. But since I don't know, this might be a safer bet.} \\
\\ 
\texttt{Wait, but the goal is to maximize the sum of both scores. So maybe the optimal split is to have each take one ruler (3 each), and then split the other items. For example, I take one ruler, glove, and one card (3+1+1=5), and the other takes the other ruler, whistle, and one card (3+1+1=5). Total combined is 10. That's the same as if I took both rulers and the other items. But in this case, the other player might have a higher} \\
            }
        }
    }
    & & \\ \\

    \theutterance \stepcounter{utterance}  
    & & \multicolumn{4}{p{0.6\linewidth}}{
        \cellcolor[rgb]{0.9,0.9,0.9}{
            \makecell[{{p{\linewidth}}}]{
                \texttt{\tiny{[P1$\langle$GM]}}
                \texttt{<think>} \\
\texttt{Okay, let's try to figure out the best way to split the items. My value function is potato: 0, glove: 1, ruler: 3, whistle: 1, card: 1. The other player's values are secret, so I need to make a proposal that's likely to be complementary.} \\
\\ 
\texttt{First, the items available are 1 potato, 1 glove, 2 rulers, 1 whistle, 2 cards. My highest value is on rulers (3 each), so I should prioritize getting as many rulers as possible. Since there are two rulers, if I take both, that's 6 points for me. Then, the remaining items are glove, whistle, and two cards. I value each of these at 1, so taking all of them would add 4 points (glove 1, whistle 1, two cards 2). But that's 6+4=10 for me. However, the other player would need to get the potato, which I value at 0. But if they value the potato, maybe they can take it. But since their values are secret, I don't know. If they also want rulers, taking both might conflict with their proposal.} \\
\\ 
\texttt{Alternatively, maybe I should take one ruler (3 points) and let the other take the other. Then, for the remaining items, I can take the glove, whistle, and one card (total 1+1+1=3), making my total 3+3=6. The other player would get the other ruler (3), the other card (1), and maybe the potato (0). Their total would be 3+1=4. Total combined is 10. But if they value the glove or whistle more, maybe they can take those. But since I don't know, this might be a safer bet.} \\
\\ 
\texttt{Wait, but the goal is to maximize the sum of both scores. So maybe the optimal split is to have each take one ruler (3 each), and then split the other items. For example, I take one ruler, glove, and one card (3+1+1=5), and the other takes the other ruler, whistle, and one card (3+1+1=5). Total combined is 10. That's the same as if I took both rulers and the other items. But in this case, the other player might have a higher} \\
            }
        }
    }
    & \\ \\

    \theutterance \stepcounter{utterance}  
    & \multicolumn{4}{p{0.6\linewidth}}{
        \cellcolor[rgb]{0.9,0.9,0.9}{
            \makecell[{{p{\linewidth}}}]{
                \texttt{\tiny{[P1$\rangle$GM]}}
                \texttt{<think>} \\
\texttt{Okay, let's try to figure out the best way to split the items. My value function is potato: 0, glove: 1, ruler: 3, whistle: 1, card: 1. The other player's values are secret, so I need to make a proposal that's likely to be complementary.} \\
\\ 
\texttt{First, the items available are 1 potato, 1 glove, 2 rulers, 1 whistle, 2 cards. My highest value is on rulers (3 each), so I should prioritize getting as many rulers as possible. Since there are two rulers, if I take both, that's 6 points for me. Then, the remaining items are glove, whistle, and two cards. I value each of these at 1, so taking all of them would add 4 points (glove 1, whistle 1, two cards 2). But that's 6+4=10 for me. However, the other player would need to get the potato, which I value at 0. But if they value the potato, maybe they can take it. But since their values are secret, I don't know. If they also want rulers, taking both might conflict with their proposal.} \\
\\ 
\texttt{Alternatively, maybe I should take one ruler (3 points) and let the other take the other. Then, for the remaining items, I can take the glove, whistle, and one card (total 1+1+1=3), making my total 3+3=6. The other player would get the other ruler (3), the other card (1), and maybe the potato (0). Their total would be 3+1=4. Total combined is 10. But if they value the glove or whistle more, maybe they can take those. But since I don't know, this might be a safer bet.} \\
\\ 
\texttt{Wait, but the goal is to maximize the sum of both scores. So maybe the optimal split is to have each take one ruler (3 each), and then split the other items. For example, I take one ruler, glove, and one card (3+1+1=5), and the other takes the other ruler, whistle, and one card (3+1+1=5). Total combined is 10. That's the same as if I took both rulers and the other items. But in this case, the other player might have a higher} \\
            }
        }
    }
    & & \\ \\

    \theutterance \stepcounter{utterance}  
    & & & \multicolumn{2}{p{0.3\linewidth}}{
        \cellcolor[rgb]{0.9,0.9,0.9}{
            \makecell[{{p{\linewidth}}}]{
                \texttt{\tiny{[GM$|$GM]}}
                \texttt{continue} \\
            }
        }
    }
    & & \\ \\

    \theutterance \stepcounter{utterance}  
    & & & \multicolumn{2}{p{0.3\linewidth}}{
        \cellcolor[rgb]{0.9,0.9,0.9}{
            \makecell[{{p{\linewidth}}}]{
                \texttt{\tiny{[GM$|$GM]}}
                \texttt{<think>} \\
\texttt{Okay, let's try to figure out the best way to split the items. My value function is potato: 0, glove: 1, ruler: 3, whistle: 1, card: 1. The other player's values are secret, so I need to make a proposal that's likely to be complementary.} \\
\\ 
\texttt{First, the items available are 1 potato, 1 glove, 2 rulers, 1 whistle, 2 cards. My highest value is on rulers (3 each), so I should prioritize getting as many rulers as possible. Since there are two rulers, if I take both, that's 6 points for me. Then, the remaining items are glove, whistle, and two cards. I value each of these at 1, so taking all of them would add 4 points (glove 1, whistle 1, two cards 2). But that's 6+4=10 for me. However, the other player would need to get the potato, which I value at 0. But if they value the potato, maybe they can take it. But since their values are secret, I don't know. If they also want rulers, taking both might conflict with their proposal.} \\
\\ 
\texttt{Alternatively, maybe I should take one ruler (3 points) and let the other take the other. Then, for the remaining items, I can take the glove, whistle, and one card (total 1+1+1=3), making my total 3+3=6. The other player would get the other ruler (3), the other card (1), and maybe the potato (0). Their total would be 3+1=4. Total combined is 10. But if they value the glove or whistle more, maybe they can take those. But since I don't know, this might be a safer bet.} \\
\\ 
\texttt{Wait, but the goal is to maximize the sum of both scores. So maybe the optimal split is to have each take one ruler (3 each), and then split the other items. For example, I take one ruler, glove, and one card (3+1+1=5), and the other takes the other ruler, whistle, and one card (3+1+1=5). Total combined is 10. That's the same as if I took both rulers and the other items. But in this case, the other player might have a higher} \\
            }
        }
    }
    & & \\ \\

    \theutterance \stepcounter{utterance}  
    & & \multicolumn{4}{p{0.6\linewidth}}{
        \cellcolor[rgb]{0.9,0.9,0.9}{
            \makecell[{{p{\linewidth}}}]{
                \texttt{\tiny{[GM$\rangle$P2]}}
                \texttt{<think>} \\
\texttt{Okay, let's try to figure out the best way to split the items. My value function is potato: 0, glove: 1, ruler: 3, whistle: 1, card: 1. The other player's values are secret, so I need to make a proposal that's likely to be complementary.} \\
\\ 
\texttt{First, the items available are 1 potato, 1 glove, 2 rulers, 1 whistle, 2 cards. My highest value is on rulers (3 each), so I should prioritize getting as many rulers as possible. Since there are two rulers, if I take both, that's 6 points for me. Then, the remaining items are glove, whistle, and two cards. I value each of these at 1, so taking all of them would add 4 points (glove 1, whistle 1, two cards 2). But that's 6+4=10 for me. However, the other player would need to get the potato, which I value at 0. But if they value the potato, maybe they can take it. But since their values are secret, I don't know. If they also want rulers, taking both might conflict with their proposal.} \\
\\ 
\texttt{Alternatively, maybe I should take one ruler (3 points) and let the other take the other. Then, for the remaining items, I can take the glove, whistle, and one card (total 1+1+1=3), making my total 3+3=6. The other player would get the other ruler (3), the other card (1), and maybe the potato (0). Their total would be 3+1=4. Total combined is 10. But if they value the glove or whistle more, maybe they can take those. But since I don't know, this might be a safer bet.} \\
\\ 
\texttt{Wait, but the goal is to maximize the sum of both scores. So maybe the optimal split is to have each take one ruler (3 each), and then split the other items. For example, I take one ruler, glove, and one card (3+1+1=5), and the other takes the other ruler, whistle, and one card (3+1+1=5). Total combined is 10. That's the same as if I took both rulers and the other items. But in this case, the other player might have a higher} \\
            }
        }
    }
    & \\ \\

    \theutterance \stepcounter{utterance}  
    & & & \multicolumn{4}{p{0.6\linewidth}}{
        \cellcolor[rgb]{0.9,0.9,0.9}{
            \makecell[{{p{\linewidth}}}]{
                \texttt{\tiny{[GM$\langle$P2]}}
                \texttt{<think>} \\
\texttt{Okay, let's try to figure out the best way to split the items. My value function is potato: 0, glove: 1, ruler: 3, whistle: 1, card: 1. The other player's values are secret, so I need to make a proposal that's likely to be complementary.} \\
\\ 
\texttt{First, the items available are 1 potato, 1 glove, 2 rulers, 1 whistle, 2 cards. My highest value is on rulers (3 each), so I should prioritize getting as many rulers as possible. Since there are two rulers, if I take both, that's 6 points for me. Then, the remaining items are glove, whistle, and two cards. I value each of these at 1, so taking all of them would add 4 points (glove 1, whistle 1, two cards 2). But that's 6+4=10 for me. However, the other player would need to get the potato, which I value at 0. But if they value the potato, maybe they can take it. But since their values are secret, I don't know. If they also want rulers, taking both might conflict with their proposal.} \\
\\ 
\texttt{Alternatively, maybe I should take one ruler (3 points) and let the other take the other. Then, for the remaining items, I can take the glove, whistle, and one card (total 1+1+1=3), making my total 3+3=6. The other player would get the other ruler (3), the other card (1), and maybe the potato (0). Their total would be 3+1=4. Total combined is 10. But if they value the glove or whistle more, maybe they can take those. But since I don't know, this might be a safer bet.} \\
\\ 
\texttt{Wait, but the goal is to maximize the sum of both scores. So maybe the optimal split is to have each take one ruler (3 each), and then split the other items. For example, I take one ruler, glove, and one card (3+1+1=5), and the other takes the other ruler, whistle, and one card (3+1+1=5). Total combined is 10. That's the same as if I took both rulers and the other items. But in this case, the other player might have a higher} \\
            }
        }
    }
     \\ \\

    \theutterance \stepcounter{utterance}  
    & & & \multicolumn{2}{p{0.3\linewidth}}{
        \cellcolor[rgb]{0.9,0.9,0.9}{
            \makecell[{{p{\linewidth}}}]{
                \texttt{\tiny{[GM$|$GM]}}
                \texttt{continue} \\
            }
        }
    }
    & & \\ \\

    \theutterance \stepcounter{utterance}  
    & & & \multicolumn{2}{p{0.3\linewidth}}{
        \cellcolor[rgb]{0.9,0.9,0.9}{
            \makecell[{{p{\linewidth}}}]{
                \texttt{\tiny{[GM$|$GM]}}
                \texttt{<think>} \\
\texttt{Okay, let's try to figure out the best way to split the items. My value function is potato: 0, glove: 1, ruler: 3, whistle: 1, card: 1. The other player's values are secret, so I need to make a proposal that's likely to be complementary.} \\
\\ 
\texttt{First, the items available are 1 potato, 1 glove, 2 rulers, 1 whistle, 2 cards. My highest value is on rulers (3 each), so I should prioritize getting as many rulers as possible. Since there are two rulers, if I take both, that's 6 points for me. Then, the remaining items are glove, whistle, and two cards. I value each of these at 1, so taking all of them would add 4 points (glove 1, whistle 1, two cards 2). But that's 6+4=10 for me. However, the other player would need to get the potato, which I value at 0. But if they value the potato, maybe they can take it. But since their values are secret, I don't know. If they also want rulers, taking both might conflict with their proposal.} \\
\\ 
\texttt{Alternatively, maybe I should take one ruler (3 points) and let the other take the other. Then, for the remaining items, I can take the glove, whistle, and one card (total 1+1+1=3), making my total 3+3=6. The other player would get the other ruler (3), the other card (1), and maybe the potato (0). Their total would be 3+1=4. Total combined is 10. But if they value the glove or whistle more, maybe they can take those. But since I don't know, this might be a safer bet.} \\
\\ 
\texttt{Wait, but the goal is to maximize the sum of both scores. So maybe the optimal split is to have each take one ruler (3 each), and then split the other items. For example, I take one ruler, glove, and one card (3+1+1=5), and the other takes the other ruler, whistle, and one card (3+1+1=5). Total combined is 10. That's the same as if I took both rulers and the other items. But in this case, the other player might have a higher} \\
            }
        }
    }
    & & \\ \\

    \theutterance \stepcounter{utterance}  
    & & \multicolumn{4}{p{0.6\linewidth}}{
        \cellcolor[rgb]{0.9,0.9,0.9}{
            \makecell[{{p{\linewidth}}}]{
                \texttt{\tiny{[P1$\langle$GM]}}
                \texttt{<think>} \\
\texttt{Okay, let's try to figure out the best way to split the items. My value function is potato: 0, glove: 1, ruler: 3, whistle: 1, card: 1. The other player's values are secret, so I need to make a proposal that's likely to be complementary.} \\
\\ 
\texttt{First, the items available are 1 potato, 1 glove, 2 rulers, 1 whistle, 2 cards. My highest value is on rulers (3 each), so I should prioritize getting as many rulers as possible. Since there are two rulers, if I take both, that's 6 points for me. Then, the remaining items are glove, whistle, and two cards. I value each of these at 1, so taking all of them would add 4 points (glove 1, whistle 1, two cards 2). But that's 6+4=10 for me. However, the other player would need to get the potato, which I value at 0. But if they value the potato, maybe they can take it. But since their values are secret, I don't know. If they also want rulers, taking both might conflict with their proposal.} \\
\\ 
\texttt{Alternatively, maybe I should take one ruler (3 points) and let the other take the other. Then, for the remaining items, I can take the glove, whistle, and one card (total 1+1+1=3), making my total 3+3=6. The other player would get the other ruler (3), the other card (1), and maybe the potato (0). Their total would be 3+1=4. Total combined is 10. But if they value the glove or whistle more, maybe they can take those. But since I don't know, this might be a safer bet.} \\
\\ 
\texttt{Wait, but the goal is to maximize the sum of both scores. So maybe the optimal split is to have each take one ruler (3 each), and then split the other items. For example, I take one ruler, glove, and one card (3+1+1=5), and the other takes the other ruler, whistle, and one card (3+1+1=5). Total combined is 10. That's the same as if I took both rulers and the other items. But in this case, the other player might have a higher} \\
            }
        }
    }
    & \\ \\

    \theutterance \stepcounter{utterance}  
    & \multicolumn{4}{p{0.6\linewidth}}{
        \cellcolor[rgb]{0.9,0.9,0.9}{
            \makecell[{{p{\linewidth}}}]{
                \texttt{\tiny{[P1$\rangle$GM]}}
                \texttt{<think>} \\
\texttt{Okay, let's try to figure out the best way to split the items. My value function is potato: 0, glove: 1, ruler: 3, whistle: 1, card: 1. The other player's values are secret, so I need to make a proposal that's likely to be complementary.} \\
\\ 
\texttt{First, the items available are 1 potato, 1 glove, 2 rulers, 1 whistle, 2 cards. My highest value is on rulers (3 each), so I should prioritize getting as many rulers as possible. Since there are two rulers, if I take both, that's 6 points for me. Then, the remaining items are glove, whistle, and two cards. I value each of these at 1, so taking all of them would add 4 points (glove 1, whistle 1, two cards 2). But that's 6+4=10 for me. However, the other player would need to get the potato, which I value at 0. But if they value the potato, maybe they can take it. But since their values are secret, I don't know. If they also want rulers, taking both might conflict with their proposal.} \\
\\ 
\texttt{Alternatively, maybe I should take one ruler (3 points) and let the other take the other. Then, for the remaining items, I can take the glove, whistle, and one card (total 1+1+1=3), making my total 3+3=6. The other player would get the other ruler (3), the other card (1), and maybe the potato (0). Their total would be 3+1=4. Total combined is 10. But if they value the glove or whistle more, maybe they can take those. But since I don't know, this might be a safer bet.} \\
\\ 
\texttt{Wait, but the goal is to maximize the sum of both scores. So maybe the optimal split is to have each take one ruler (3 each), and then split the other items. For example, I take one ruler, glove, and one card (3+1+1=5), and the other takes the other ruler, whistle, and one card (3+1+1=5). Total combined is 10. That's the same as if I took both rulers and the other items. But in this case, the other player might have a higher} \\
            }
        }
    }
    & & \\ \\

    \theutterance \stepcounter{utterance}  
    & & & \multicolumn{2}{p{0.3\linewidth}}{
        \cellcolor[rgb]{0.9,0.9,0.9}{
            \makecell[{{p{\linewidth}}}]{
                \texttt{\tiny{[GM$|$GM]}}
                \texttt{continue} \\
            }
        }
    }
    & & \\ \\

    \theutterance \stepcounter{utterance}  
    & & & \multicolumn{2}{p{0.3\linewidth}}{
        \cellcolor[rgb]{0.9,0.9,0.9}{
            \makecell[{{p{\linewidth}}}]{
                \texttt{\tiny{[GM$|$GM]}}
                \texttt{<think>} \\
\texttt{Okay, let's try to figure out the best way to split the items. My value function is potato: 0, glove: 1, ruler: 3, whistle: 1, card: 1. The other player's values are secret, so I need to make a proposal that's likely to be complementary.} \\
\\ 
\texttt{First, the items available are 1 potato, 1 glove, 2 rulers, 1 whistle, 2 cards. My highest value is on rulers (3 each), so I should prioritize getting as many rulers as possible. Since there are two rulers, if I take both, that's 6 points for me. Then, the remaining items are glove, whistle, and two cards. I value each of these at 1, so taking all of them would add 4 points (glove 1, whistle 1, two cards 2). But that's 6+4=10 for me. However, the other player would need to get the potato, which I value at 0. But if they value the potato, maybe they can take it. But since their values are secret, I don't know. If they also want rulers, taking both might conflict with their proposal.} \\
\\ 
\texttt{Alternatively, maybe I should take one ruler (3 points) and let the other take the other. Then, for the remaining items, I can take the glove, whistle, and one card (total 1+1+1=3), making my total 3+3=6. The other player would get the other ruler (3), the other card (1), and maybe the potato (0). Their total would be 3+1=4. Total combined is 10. But if they value the glove or whistle more, maybe they can take those. But since I don't know, this might be a safer bet.} \\
\\ 
\texttt{Wait, but the goal is to maximize the sum of both scores. So maybe the optimal split is to have each take one ruler (3 each), and then split the other items. For example, I take one ruler, glove, and one card (3+1+1=5), and the other takes the other ruler, whistle, and one card (3+1+1=5). Total combined is 10. That's the same as if I took both rulers and the other items. But in this case, the other player might have a higher} \\
            }
        }
    }
    & & \\ \\

    \theutterance \stepcounter{utterance}  
    & & \multicolumn{4}{p{0.6\linewidth}}{
        \cellcolor[rgb]{0.9,0.9,0.9}{
            \makecell[{{p{\linewidth}}}]{
                \texttt{\tiny{[GM$\rangle$P2]}}
                \texttt{<think>} \\
\texttt{Okay, let's try to figure out the best way to split the items. My value function is potato: 0, glove: 1, ruler: 3, whistle: 1, card: 1. The other player's values are secret, so I need to make a proposal that's likely to be complementary.} \\
\\ 
\texttt{First, the items available are 1 potato, 1 glove, 2 rulers, 1 whistle, 2 cards. My highest value is on rulers (3 each), so I should prioritize getting as many rulers as possible. Since there are two rulers, if I take both, that's 6 points for me. Then, the remaining items are glove, whistle, and two cards. I value each of these at 1, so taking all of them would add 4 points (glove 1, whistle 1, two cards 2). But that's 6+4=10 for me. However, the other player would need to get the potato, which I value at 0. But if they value the potato, maybe they can take it. But since their values are secret, I don't know. If they also want rulers, taking both might conflict with their proposal.} \\
\\ 
\texttt{Alternatively, maybe I should take one ruler (3 points) and let the other take the other. Then, for the remaining items, I can take the glove, whistle, and one card (total 1+1+1=3), making my total 3+3=6. The other player would get the other ruler (3), the other card (1), and maybe the potato (0). Their total would be 3+1=4. Total combined is 10. But if they value the glove or whistle more, maybe they can take those. But since I don't know, this might be a safer bet.} \\
\\ 
\texttt{Wait, but the goal is to maximize the sum of both scores. So maybe the optimal split is to have each take one ruler (3 each), and then split the other items. For example, I take one ruler, glove, and one card (3+1+1=5), and the other takes the other ruler, whistle, and one card (3+1+1=5). Total combined is 10. That's the same as if I took both rulers and the other items. But in this case, the other player might have a higher} \\
            }
        }
    }
    & \\ \\

    \theutterance \stepcounter{utterance}  
    & & & \multicolumn{4}{p{0.6\linewidth}}{
        \cellcolor[rgb]{0.9,0.9,0.9}{
            \makecell[{{p{\linewidth}}}]{
                \texttt{\tiny{[GM$\langle$P2]}}
                \texttt{<think>} \\
\texttt{Okay, let's try to figure out the best way to split the items. My value function is potato: 0, glove: 1, ruler: 3, whistle: 1, card: 1. The other player's values are secret, so I need to make a proposal that's likely to be complementary.} \\
\\ 
\texttt{First, the items available are 1 potato, 1 glove, 2 rulers, 1 whistle, 2 cards. My highest value is on rulers (3 each), so I should prioritize getting as many rulers as possible. Since there are two rulers, if I take both, that's 6 points for me. Then, the remaining items are glove, whistle, and two cards. I value each of these at 1, so taking all of them would add 4 points (glove 1, whistle 1, two cards 2). But that's 6+4=10 for me. However, the other player would need to get the potato, which I value at 0. But if they value the potato, maybe they can take it. But since their values are secret, I don't know. If they also want rulers, taking both might conflict with their proposal.} \\
\\ 
\texttt{Alternatively, maybe I should take one ruler (3 points) and let the other take the other. Then, for the remaining items, I can take the glove, whistle, and one card (total 1+1+1=3), making my total 3+3=6. The other player would get the other ruler (3), the other card (1), and maybe the potato (0). Their total would be 3+1=4. Total combined is 10. But if they value the glove or whistle more, maybe they can take those. But since I don't know, this might be a safer bet.} \\
\\ 
\texttt{Wait, but the goal is to maximize the sum of both scores. So maybe the optimal split is to have each take one ruler (3 each), and then split the other items. For example, I take one ruler, glove, and one card (3+1+1=5), and the other takes the other ruler, whistle, and one card (3+1+1=5). Total combined is 10. That's the same as if I took both rulers and the other items. But in this case, the other player might have a higher} \\
            }
        }
    }
     \\ \\

    \theutterance \stepcounter{utterance}  
    & & & \multicolumn{2}{p{0.3\linewidth}}{
        \cellcolor[rgb]{0.9,0.9,0.9}{
            \makecell[{{p{\linewidth}}}]{
                \texttt{\tiny{[GM$|$GM]}}
                \texttt{continue} \\
            }
        }
    }
    & & \\ \\

    \theutterance \stepcounter{utterance}  
    & & & \multicolumn{2}{p{0.3\linewidth}}{
        \cellcolor[rgb]{0.9,0.9,0.9}{
            \makecell[{{p{\linewidth}}}]{
                \texttt{\tiny{[GM$|$GM]}}
                \texttt{<think>} \\
\texttt{Okay, let's try to figure out the best way to split the items. My value function is potato: 0, glove: 1, ruler: 3, whistle: 1, card: 1. The other player's values are secret, so I need to make a proposal that's likely to be complementary.} \\
\\ 
\texttt{First, the items available are 1 potato, 1 glove, 2 rulers, 1 whistle, 2 cards. My highest value is on rulers (3 each), so I should prioritize getting as many rulers as possible. Since there are two rulers, if I take both, that's 6 points for me. Then, the remaining items are glove, whistle, and two cards. I value each of these at 1, so taking all of them would add 4 points (glove 1, whistle 1, two cards 2). But that's 6+4=10 for me. However, the other player would need to get the potato, which I value at 0. But if they value the potato, maybe they can take it. But since their values are secret, I don't know. If they also want rulers, taking both might conflict with their proposal.} \\
\\ 
\texttt{Alternatively, maybe I should take one ruler (3 points) and let the other take the other. Then, for the remaining items, I can take the glove, whistle, and one card (total 1+1+1=3), making my total 3+3=6. The other player would get the other ruler (3), the other card (1), and maybe the potato (0). Their total would be 3+1=4. Total combined is 10. But if they value the glove or whistle more, maybe they can take those. But since I don't know, this might be a safer bet.} \\
\\ 
\texttt{Wait, but the goal is to maximize the sum of both scores. So maybe the optimal split is to have each take one ruler (3 each), and then split the other items. For example, I take one ruler, glove, and one card (3+1+1=5), and the other takes the other ruler, whistle, and one card (3+1+1=5). Total combined is 10. That's the same as if I took both rulers and the other items. But in this case, the other player might have a higher} \\
            }
        }
    }
    & & \\ \\

    \theutterance \stepcounter{utterance}  
    & & \multicolumn{4}{p{0.6\linewidth}}{
        \cellcolor[rgb]{0.9,0.9,0.9}{
            \makecell[{{p{\linewidth}}}]{
                \texttt{\tiny{[P1$\langle$GM]}}
                \texttt{<think>} \\
\texttt{Okay, let's try to figure out the best way to split the items. My value function is potato: 0, glove: 1, ruler: 3, whistle: 1, card: 1. The other player's values are secret, so I need to make a proposal that's likely to be complementary.} \\
\\ 
\texttt{First, the items available are 1 potato, 1 glove, 2 rulers, 1 whistle, 2 cards. My highest value is on rulers (3 each), so I should prioritize getting as many rulers as possible. Since there are two rulers, if I take both, that's 6 points for me. Then, the remaining items are glove, whistle, and two cards. I value each of these at 1, so taking all of them would add 4 points (glove 1, whistle 1, two cards 2). But that's 6+4=10 for me. However, the other player would need to get the potato, which I value at 0. But if they value the potato, maybe they can take it. But since their values are secret, I don't know. If they also want rulers, taking both might conflict with their proposal.} \\
\\ 
\texttt{Alternatively, maybe I should take one ruler (3 points) and let the other take the other. Then, for the remaining items, I can take the glove, whistle, and one card (total 1+1+1=3), making my total 3+3=6. The other player would get the other ruler (3), the other card (1), and maybe the potato (0). Their total would be 3+1=4. Total combined is 10. But if they value the glove or whistle more, maybe they can take those. But since I don't know, this might be a safer bet.} \\
\\ 
\texttt{Wait, but the goal is to maximize the sum of both scores. So maybe the optimal split is to have each take one ruler (3 each), and then split the other items. For example, I take one ruler, glove, and one card (3+1+1=5), and the other takes the other ruler, whistle, and one card (3+1+1=5). Total combined is 10. That's the same as if I took both rulers and the other items. But in this case, the other player might have a higher} \\
            }
        }
    }
    & \\ \\

    \theutterance \stepcounter{utterance}  
    & \multicolumn{4}{p{0.6\linewidth}}{
        \cellcolor[rgb]{0.9,0.9,0.9}{
            \makecell[{{p{\linewidth}}}]{
                \texttt{\tiny{[P1$\rangle$GM]}}
                \texttt{<think>} \\
\texttt{Okay, let's try to figure out the best way to split the items. My value function is potato: 0, glove: 1, ruler: 3, whistle: 1, card: 1. The other player's values are secret, so I need to make a proposal that's likely to be complementary.} \\
\\ 
\texttt{First, the items available are 1 potato, 1 glove, 2 rulers, 1 whistle, 2 cards. My highest value is on rulers (3 each), so I should prioritize getting as many rulers as possible. Since there are two rulers, if I take both, that's 6 points for me. Then, the remaining items are glove, whistle, and two cards. I value each of these at 1, so taking all of them would add 4 points (glove 1, whistle 1, two cards 2). But that's 6+4=10 for me. However, the other player would need to get the potato, which I value at 0. But if they value the potato, maybe they can take it. But since their values are secret, I don't know. If they also want rulers, taking both might conflict with their proposal.} \\
\\ 
\texttt{Alternatively, maybe I should take one ruler (3 points) and let the other take the other. Then, for the remaining items, I can take the glove, whistle, and one card (total 1+1+1=3), making my total 3+3=6. The other player would get the other ruler (3), the other card (1), and maybe the potato (0). Their total would be 3+1=4. Total combined is 10. But if they value the glove or whistle more, maybe they can take those. But since I don't know, this might be a safer bet.} \\
\\ 
\texttt{Wait, but the goal is to maximize the sum of both scores. So maybe the optimal split is to have each take one ruler (3 each), and then split the other items. For example, I take one ruler, glove, and one card (3+1+1=5), and the other takes the other ruler, whistle, and one card (3+1+1=5). Total combined is 10. That's the same as if I took both rulers and the other items. But in this case, the other player might have a higher} \\
            }
        }
    }
    & & \\ \\

    \theutterance \stepcounter{utterance}  
    & & & \multicolumn{2}{p{0.3\linewidth}}{
        \cellcolor[rgb]{0.9,0.9,0.9}{
            \makecell[{{p{\linewidth}}}]{
                \texttt{\tiny{[GM$|$GM]}}
                \texttt{continue} \\
            }
        }
    }
    & & \\ \\

    \theutterance \stepcounter{utterance}  
    & & & \multicolumn{2}{p{0.3\linewidth}}{
        \cellcolor[rgb]{0.9,0.9,0.9}{
            \makecell[{{p{\linewidth}}}]{
                \texttt{\tiny{[GM$|$GM]}}
                \texttt{<think>} \\
\texttt{Okay, let's try to figure out the best way to split the items. My value function is potato: 0, glove: 1, ruler: 3, whistle: 1, card: 1. The other player's values are secret, so I need to make a proposal that's likely to be complementary.} \\
\\ 
\texttt{First, the items available are 1 potato, 1 glove, 2 rulers, 1 whistle, 2 cards. My highest value is on rulers (3 each), so I should prioritize getting as many rulers as possible. Since there are two rulers, if I take both, that's 6 points for me. Then, the remaining items are glove, whistle, and two cards. I value each of these at 1, so taking all of them would add 4 points (glove 1, whistle 1, two cards 2). But that's 6+4=10 for me. However, the other player would need to get the potato, which I value at 0. But if they value the potato, maybe they can take it. But since their values are secret, I don't know. If they also want rulers, taking both might conflict with their proposal.} \\
\\ 
\texttt{Alternatively, maybe I should take one ruler (3 points) and let the other take the other. Then, for the remaining items, I can take the glove, whistle, and one card (total 1+1+1=3), making my total 3+3=6. The other player would get the other ruler (3), the other card (1), and maybe the potato (0). Their total would be 3+1=4. Total combined is 10. But if they value the glove or whistle more, maybe they can take those. But since I don't know, this might be a safer bet.} \\
\\ 
\texttt{Wait, but the goal is to maximize the sum of both scores. So maybe the optimal split is to have each take one ruler (3 each), and then split the other items. For example, I take one ruler, glove, and one card (3+1+1=5), and the other takes the other ruler, whistle, and one card (3+1+1=5). Total combined is 10. That's the same as if I took both rulers and the other items. But in this case, the other player might have a higher} \\
            }
        }
    }
    & & \\ \\

    \theutterance \stepcounter{utterance}  
    & & \multicolumn{4}{p{0.6\linewidth}}{
        \cellcolor[rgb]{0.9,0.9,0.9}{
            \makecell[{{p{\linewidth}}}]{
                \texttt{\tiny{[GM$\rangle$P2]}}
                \texttt{<think>} \\
\texttt{Okay, let's try to figure out the best way to split the items. My value function is potato: 0, glove: 1, ruler: 3, whistle: 1, card: 1. The other player's values are secret, so I need to make a proposal that's likely to be complementary.} \\
\\ 
\texttt{First, the items available are 1 potato, 1 glove, 2 rulers, 1 whistle, 2 cards. My highest value is on rulers (3 each), so I should prioritize getting as many rulers as possible. Since there are two rulers, if I take both, that's 6 points for me. Then, the remaining items are glove, whistle, and two cards. I value each of these at 1, so taking all of them would add 4 points (glove 1, whistle 1, two cards 2). But that's 6+4=10 for me. However, the other player would need to get the potato, which I value at 0. But if they value the potato, maybe they can take it. But since their values are secret, I don't know. If they also want rulers, taking both might conflict with their proposal.} \\
\\ 
\texttt{Alternatively, maybe I should take one ruler (3 points) and let the other take the other. Then, for the remaining items, I can take the glove, whistle, and one card (total 1+1+1=3), making my total 3+3=6. The other player would get the other ruler (3), the other card (1), and maybe the potato (0). Their total would be 3+1=4. Total combined is 10. But if they value the glove or whistle more, maybe they can take those. But since I don't know, this might be a safer bet.} \\
\\ 
\texttt{Wait, but the goal is to maximize the sum of both scores. So maybe the optimal split is to have each take one ruler (3 each), and then split the other items. For example, I take one ruler, glove, and one card (3+1+1=5), and the other takes the other ruler, whistle, and one card (3+1+1=5). Total combined is 10. That's the same as if I took both rulers and the other items. But in this case, the other player might have a higher} \\
            }
        }
    }
    & \\ \\

    \theutterance \stepcounter{utterance}  
    & & & \multicolumn{4}{p{0.6\linewidth}}{
        \cellcolor[rgb]{0.9,0.9,0.9}{
            \makecell[{{p{\linewidth}}}]{
                \texttt{\tiny{[GM$\langle$P2]}}
                \texttt{<think>} \\
\texttt{Okay, let's try to figure out the best way to split the items. My value function is potato: 0, glove: 1, ruler: 3, whistle: 1, card: 1. The other player's values are secret, so I need to make a proposal that's likely to be complementary.} \\
\\ 
\texttt{First, the items available are 1 potato, 1 glove, 2 rulers, 1 whistle, 2 cards. My highest value is on rulers (3 each), so I should prioritize getting as many rulers as possible. Since there are two rulers, if I take both, that's 6 points for me. Then, the remaining items are glove, whistle, and two cards. I value each of these at 1, so taking all of them would add 4 points (glove 1, whistle 1, two cards 2). But that's 6+4=10 for me. However, the other player would need to get the potato, which I value at 0. But if they value the potato, maybe they can take it. But since their values are secret, I don't know. If they also want rulers, taking both might conflict with their proposal.} \\
\\ 
\texttt{Alternatively, maybe I should take one ruler (3 points) and let the other take the other. Then, for the remaining items, I can take the glove, whistle, and one card (total 1+1+1=3), making my total 3+3=6. The other player would get the other ruler (3), the other card (1), and maybe the potato (0). Their total would be 3+1=4. Total combined is 10. But if they value the glove or whistle more, maybe they can take those. But since I don't know, this might be a safer bet.} \\
\\ 
\texttt{Wait, but the goal is to maximize the sum of both scores. So maybe the optimal split is to have each take one ruler (3 each), and then split the other items. For example, I take one ruler, glove, and one card (3+1+1=5), and the other takes the other ruler, whistle, and one card (3+1+1=5). Total combined is 10. That's the same as if I took both rulers and the other items. But in this case, the other player might have a higher} \\
            }
        }
    }
     \\ \\

    \theutterance \stepcounter{utterance}  
    & & & \multicolumn{2}{p{0.3\linewidth}}{
        \cellcolor[rgb]{0.9,0.9,0.9}{
            \makecell[{{p{\linewidth}}}]{
                \texttt{\tiny{[GM$|$GM]}}
                \texttt{continue} \\
            }
        }
    }
    & & \\ \\

    \theutterance \stepcounter{utterance}  
    & & & \multicolumn{2}{p{0.3\linewidth}}{
        \cellcolor[rgb]{0.9,0.9,0.9}{
            \makecell[{{p{\linewidth}}}]{
                \texttt{\tiny{[GM$|$GM]}}
                \texttt{<think>} \\
\texttt{Okay, let's try to figure out the best way to split the items. My value function is potato: 0, glove: 1, ruler: 3, whistle: 1, card: 1. The other player's values are secret, so I need to make a proposal that's likely to be complementary.} \\
\\ 
\texttt{First, the items available are 1 potato, 1 glove, 2 rulers, 1 whistle, 2 cards. My highest value is on rulers (3 each), so I should prioritize getting as many rulers as possible. Since there are two rulers, if I take both, that's 6 points for me. Then, the remaining items are glove, whistle, and two cards. I value each of these at 1, so taking all of them would add 4 points (glove 1, whistle 1, two cards 2). But that's 6+4=10 for me. However, the other player would need to get the potato, which I value at 0. But if they value the potato, maybe they can take it. But since their values are secret, I don't know. If they also want rulers, taking both might conflict with their proposal.} \\
\\ 
\texttt{Alternatively, maybe I should take one ruler (3 points) and let the other take the other. Then, for the remaining items, I can take the glove, whistle, and one card (total 1+1+1=3), making my total 3+3=6. The other player would get the other ruler (3), the other card (1), and maybe the potato (0). Their total would be 3+1=4. Total combined is 10. But if they value the glove or whistle more, maybe they can take those. But since I don't know, this might be a safer bet.} \\
\\ 
\texttt{Wait, but the goal is to maximize the sum of both scores. So maybe the optimal split is to have each take one ruler (3 each), and then split the other items. For example, I take one ruler, glove, and one card (3+1+1=5), and the other takes the other ruler, whistle, and one card (3+1+1=5). Total combined is 10. That's the same as if I took both rulers and the other items. But in this case, the other player might have a higher} \\
            }
        }
    }
    & & \\ \\

    \theutterance \stepcounter{utterance}  
    & & \multicolumn{4}{p{0.6\linewidth}}{
        \cellcolor[rgb]{0.9,0.9,0.9}{
            \makecell[{{p{\linewidth}}}]{
                \texttt{\tiny{[P1$\langle$GM]}}
                \texttt{<think>} \\
\texttt{Okay, let's try to figure out the best way to split the items. My value function is potato: 0, glove: 1, ruler: 3, whistle: 1, card: 1. The other player's values are secret, so I need to make a proposal that's likely to be complementary.} \\
\\ 
\texttt{First, the items available are 1 potato, 1 glove, 2 rulers, 1 whistle, 2 cards. My highest value is on rulers (3 each), so I should prioritize getting as many rulers as possible. Since there are two rulers, if I take both, that's 6 points for me. Then, the remaining items are glove, whistle, and two cards. I value each of these at 1, so taking all of them would add 4 points (glove 1, whistle 1, two cards 2). But that's 6+4=10 for me. However, the other player would need to get the potato, which I value at 0. But if they value the potato, maybe they can take it. But since their values are secret, I don't know. If they also want rulers, taking both might conflict with their proposal.} \\
\\ 
\texttt{Alternatively, maybe I should take one ruler (3 points) and let the other take the other. Then, for the remaining items, I can take the glove, whistle, and one card (total 1+1+1=3), making my total 3+3=6. The other player would get the other ruler (3), the other card (1), and maybe the potato (0). Their total would be 3+1=4. Total combined is 10. But if they value the glove or whistle more, maybe they can take those. But since I don't know, this might be a safer bet.} \\
\\ 
\texttt{Wait, but the goal is to maximize the sum of both scores. So maybe the optimal split is to have each take one ruler (3 each), and then split the other items. For example, I take one ruler, glove, and one card (3+1+1=5), and the other takes the other ruler, whistle, and one card (3+1+1=5). Total combined is 10. That's the same as if I took both rulers and the other items. But in this case, the other player might have a higher} \\
\\ 
\\ 
\texttt{The time is up. It is now your turn to submit a secret proposal.} \\
            }
        }
    }
    & \\ \\

    \theutterance \stepcounter{utterance}  
    & \multicolumn{4}{p{0.6\linewidth}}{
        \cellcolor[rgb]{0.9,0.9,0.9}{
            \makecell[{{p{\linewidth}}}]{
                \texttt{\tiny{[P1$\rangle$GM]}}
                \texttt{<think>} \\
\texttt{Okay, let's try to figure out the best way to split the items. My value function is potato: 0, glove: 1, ruler: 3, whistle: 1, card: 1. The other player's values are secret, so I need to make a proposal that's likely to be complementary.} \\
\\ 
\texttt{First, the items available are 1 potato, 1 glove, 2 rulers, 1 whistle, 2 cards. My highest value is on rulers (3 each), so I should prioritize getting as many rulers as possible. Since there are two rulers, if I take both, that's 6 points for me. Then, the remaining items are glove, whistle, and two cards. I value each of these at 1, so taking all of them would add 4 points (glove 1, whistle 1, two cards 2). But that's 6+4=10 for me. However, the other player would need to get the potato, which I value at 0. But if they value the potato, maybe they can take it. But since their values are secret, I don't know. If they also want rulers, taking both might conflict with their proposal.} \\
\\ 
\texttt{Alternatively, maybe I should take one ruler (3 points) and let the other take the other. Then, for the remaining items, I can take the glove, whistle, and one card (total 1+1+1=3), making my total 3+3=6. The other player would get the other ruler (3), the other card (1), and maybe the potato (0). Their total would be 3+1=4. Total combined is 10. But if they value the glove or whistle more, maybe they can take those. But since I don't know, this might be a safer bet.} \\
\\ 
\texttt{Wait, but the goal is to maximize the sum of both scores. So maybe the optimal split is to have each take one ruler (3 each), and then split the other items. For example, I take one ruler, glove, and one card (3+1+1=5), and the other takes the other ruler, whistle, and one card (3+1+1=5). Total combined is 10. That's the same as if I took both rulers and the other items. But in this case, the other player might have a higher} \\
            }
        }
    }
    & & \\ \\

    \theutterance \stepcounter{utterance}  
    & & & \multicolumn{2}{p{0.3\linewidth}}{
        \cellcolor[rgb]{0.9,0.9,0.9}{
            \makecell[{{p{\linewidth}}}]{
                \texttt{\tiny{[GM$|$GM]}}
                \texttt{continue} \\
            }
        }
    }
    & & \\ \\

    \theutterance \stepcounter{utterance}  
    & & & \multicolumn{2}{p{0.3\linewidth}}{
        \cellcolor[rgb]{0.9,0.9,0.9}{
            \makecell[{{p{\linewidth}}}]{
                \texttt{\tiny{[GM$|$GM]}}
                \texttt{player was instructed to make a proposal} \\
            }
        }
    }
    & & \\ \\

\end{supertabular}
}

\end{document}
